% 分野別類題: 一次同次線形微分方程式 解答例
% コンパイル: TEXINPUTS="../../../00_common:$TEXINPUTS" latexmk -xelatex answers.tex
\documentclass[a4paper,11pt]{article}
\usepackage{preamble}
\usepackage{shoakubox}

\begin{document}

\kamokumei{類題演習 解答例 --- 一次同次線形微分方程式}

\daimon{【解法】一次同次線形微分方程式の解き方と導出}

\noindent
一次同次線形微分方程式とは、次の形の1階微分方程式である。
\[
\frac{dy}{dx} + P(x)\,y = 0
\]
変数分離形であるから、$y \neq 0$ のとき両辺を $y$ で割って
\[
\frac{dy}{y} = -P(x)\,dx
\]
両辺を積分すると
\[
\ln|y| = -\int P(x)\,dx + C_{1}
\quad\Longrightarrow\quad
y = C\exp\left(-\int P(x)\,dx\right)
\]
（$C = \pm e^{C_{1}}$。$C = 0$ とすれば自明解 $y = 0$ も含む。）初期条件が与えられた場合は、一般解に代入して $C$ を定める。

\clearpage
\begin{mondaiboxS}{問1}
$\dfrac{dy}{dx} + 3x^{2}\,y = 0$ の一般解を求めよ。
\end{mondaiboxS}
\kaitou
変数分離して（$y \neq 0$）
\[
\frac{dy}{y} = -3x^{2}\,dx
\quad\Longrightarrow\quad
\ln|y| = -x^{3} + C_{1}
\]
よって
\[
\boxed{\; y = Ce^{-x^{3}} \;}
\qquad (C \text{ は任意定数。} C = 0 \text{ とすれば } y = 0 \text{ も含む})
\]

\clearpage
\begin{mondaiboxS}{問2}
$x\dfrac{dy}{dx} + 2y = 0 \quad (x > 0)$ の一般解を求めよ。
\end{mondaiboxS}
\kaitou
$\dfrac{dy}{dx} + \dfrac{2}{x}\,y = 0$ と変形し、変数分離して
\[
\frac{dy}{y} = -\frac{2}{x}\,dx
\quad\Longrightarrow\quad
\ln|y| = -2\ln x + C_{1}
\]
よって
\[
\boxed{\; y = \frac{C}{x^{2}} \;}
\qquad (C \text{ は任意定数})
\]

\clearpage
\begin{mondaiboxS}{問3}
$\dfrac{dy}{dx} - y\tan x = 0, \quad y(0) = 2 \quad \left(-\dfrac{\pi}{2} < x < \dfrac{\pi}{2}\right)$ を解け。
\end{mondaiboxS}
\kaitou
変数分離して
\[
\frac{dy}{y} = \tan x\,dx
\quad\Longrightarrow\quad
\ln|y| = -\ln|\cos x| + C_{1}
\]
よって一般解は $y = \dfrac{C}{\cos x}$。初期条件 $y(0) = 2$ より $C = 2$。したがって
\[
\boxed{\; y = \frac{2}{\cos x} \;}
\qquad \left(-\frac{\pi}{2} < x < \frac{\pi}{2}\right)
\]

\end{document}
